%LTeX: language=fr-FR
\documentclass[12pt]{article}
\usepackage[margin=0.8in]{geometry}
\usepackage[utf8]{inputenc}
\usepackage{graphicx}
\usepackage[export]{adjustbox}
\usepackage{amsmath}
\usepackage{amsfonts}
\usepackage{amssymb}
\usepackage[version=4]{mhchem}
\usepackage{stmaryrd}
\usepackage{bbold}
%\usepackage[fallback]{xeCJK}
%\usepackage{polyglossia}
\usepackage{fontspec}
\usepackage{censor}
\usepackage{currfile}
%\setCJKmainfont{Noto Serif CJK TC}
\usepackage{enumitem}
\usepackage{tikz}
\usetikzlibrary {arrows.meta, quotes, calc} 
\usepackage{pgfplots}
\pgfplotsset{compat=1.18}
\usepackage{amsthm}
\usepackage{tcolorbox}
\usepackage{array}
\usepackage{tabularray}
\usepackage{makecell}
\usepackage{esvect}
%\usepackage{changepage}   % for the adjustwidth environment
\usepackage{enumitem}
\usepackage{tasks}
\usepackage{fancyhdr}
\usepackage[french]{babel}
\usepackage{lastpage}
\usepackage{needspace}

\tcbuselibrary{theorems}

\definecolor{GrayCen}{HTML}{d2d2d2}
\definecolor{GrayBG}{HTML}{d9d9d9}
\definecolor{BlueDef}{HTML}{104e8b}
\definecolor{GreenProp}{HTML}{025540}
\definecolor{RedMet}{HTML}{cd5556}
\definecolor{BlackSav}{HTML}{000000}
\def\censorcolor{GrayCen}
\let\svcensorrule\censorrule
\renewcommand\censorrule[1]{%
\textcolor{\censorcolor}{\svcensorrule{#1}}}

\newtcbtheorem
[]% init options
{definition}% name
{Définition}% title
{%
  colback=GrayBG,
  colframe=BlueDef,
  fonttitle=\bfseries,
}% options
{def}% prefix

\newtcbtheorem
[]% init options
{proposition}% name
{Proposition}% title
{%
  colback=GrayBG,
  colframe=GreenProp,
  fonttitle=\bfseries,
}% options
{prop}% prefix

\newtcbtheorem
[]% init options
{method}% name
{Méthode}% title
{%
  theorem name,
  colback=GrayBG,
  colframe=RedMet,
  fonttitle=\bfseries,
}% options
{met}% prefix

\newtcbtheorem
[]% init options
{savoir}% name
{Savoir-faire du chapitre}% title
{%
  theorem name,
  colback=GrayBG,
  colframe=BlackSav,
  fonttitle=\bfseries,
}% options
{sav}% prefix

\newtheorem*{exampleil}{Exemple}

\newtheoremstyle{break}
{\topsep}{\topsep}%
{}{}%
{\bfseries}{}%
{\newline}{}%
\theoremstyle{break}
\newtheorem*{example}{Exemple.}
\newtheorem*{examples}{Exemples.}
\newtheorem*{remark}{Remarque.}
\newtheorem*{remarks}{Remarques.}

\newenvironment{demo}[1][\proofname]{%
  \begin{proof}[#1]$ $\newline
  }{%
  \end{proof}
}

%\setmainfont{CMU Serif}name

\title{
  \vspace{4em}
  \centering
  \tcbset{colframe=BlackSav,colback=GrayBG}
  \tcbox[tikznode]{
    Chapitre 2 \\Les ensembles de nombres
  }
}

\author{}
\date{}

\pagestyle{empty}
\graphicspath{ {../Images/} }
\input{\currfiledir ../def.tex}
\def \Document {Cours}

\setlength\parindent{0pt}
%\setlist[itemize]{topsep=0pt}
%\setlist{nolistsep}
%\setlist[itemize]{noitemsep}
%\setlist[enumerate]{noitemsep}
\setlist[itemize,1]{label=$\bullet$}

\tikzset{cross/.style={path picture={
  \draw[black]
    (path picture bounding box.south east)--(path picture bounding box.north west)
    (path picture bounding box.south west)--(path picture bounding box.north east);
}}}

\newcolumntype{C}[1]{>{\centering\arraybackslash}m{#1}}
\newcolumntype{L}[1]{>{\raggedright\arraybackslash}m{#1}}
\newcolumntype{N}{@{}m{0pt}@{}}

\fancypagestyle{firststyle}
{
  \fancyhf{}
  \fancyhead[L]{\Structure}
  \fancyhead[R]{\Niveau}
  \fancyfoot[L]{Alexandre Fernandez}
  \fancyfoot[R]{{\thepage\  / \pageref{LastPage}}}
   %\renewcommand{\headrulewidth}{0pt} % removes horizontal header line
}


\begin{document}

\maketitle
\thispagestyle{firststyle}

\tableofcontents

% On crée une boîte temporaire pour stocker le texte
\newsavebox{\hidebox}

% -----------
% Définition de l'environnement qui rend le contenu invisible
% \newcommand{\hideM}[1]{\mbox{\vphantom{$#1$}\smash{\colorbox{gray!50}{\phantom{$\displaystyle #1$}}}} }
% \newcommand{\hideMW}[1]{\mbox{\vphantom{$#1$}\smash{\colorbox{white}{\phantom{$\displaystyle #1$}}}} }
% \newcommand{\hide}[1]{\mbox{\vphantom{#1}\smash{\colorbox{gray!50}{\phantom{#1}}}}}
% \newcommand{\hideW}[1]{\mbox{\vphantom{#1}\smash{\colorbox{white}{\phantom{#1}}}}}
% \NewEnviron{multihide}{%
%   \colorbox{gray!50}{
%     \begin{lrbox}{\hidebox}%
%         \begin{minipage}{\linewidth}%
%           \BODY
%         \end{minipage}%
%     \end{lrbox}%
%     % On affiche un espace invisible ayant les dimensions exactes de la boîte
%     \vrule width 0pt height \ht\hidebox depth \dp\hidebox\hskip\wd\hidebox%
%   }
% }
% \NewEnviron{multihideW}{%
%   \colorbox{white}{
%     \begin{lrbox}{\hidebox}%
%         \begin{minipage}{\linewidth}%
%           \BODY
%         \end{minipage}%
%     \end{lrbox}%
%     % On affiche un espace invisible ayant les dimensions exactes de la boîte
%     \vrule width 0pt height \ht\hidebox depth \dp\hidebox\hskip\wd\hidebox%
%   }
% }
% \newcommand{\hidetk}[1]{}
% Définition de l'environnement qui rend le contenu invisible
\newcommand{\hideM}[1]{\mbox{\vphantom{$#1$}\smash{\colorbox{gray!50}{$\displaystyle #1$}}} }
\newcommand{\hideMW}[1]{\mbox{\vphantom{$#1$}\smash{\colorbox{white}{$\displaystyle #1$}}} }
\newcommand{\hide}[1]{\mbox{\vphantom{#1}\smash{\colorbox{gray!50}{#1}}}}
\newcommand{\hideW}[1]{\mbox{\vphantom{#1}\smash{#1}}}
\NewEnviron{multihide}{%
  \colorbox{gray!50}{
    \begin{minipage}{\linewidth}%
      \BODY
    \end{minipage}%
  }
}
\NewEnviron{multihideW}{%
  \colorbox{white}{
    \begin{minipage}{\linewidth}%
      \BODY
    \end{minipage}%
  }
}
\newcommand{\hidetk}[1]{#1}
% -----------

\newcommand{\wvec}[1]{\overrightarrow{#1}}
\newcommand{\ssi}[0]{\emph{ssi}\ }

% Macros pour la représentation des intervalles sur la droite réelle
% (crochet gauche fermé / ouvert, crochet droit fermé / ouvert)
\newcommand{\droite}{\draw[<->, >=Stealth] (0,0) -- (6,0);}
\newcommand{\brGF}[1]{\node[inner sep=0pt] at (#1,0) {\Large$[$};}
\newcommand{\brGO}[1]{\node[inner sep=0pt] at (#1,0) {\Large$]$};}
\newcommand{\brDF}[1]{\node[inner sep=0pt] at (#1,0) {\Large$]$};}
\newcommand{\brDO}[1]{\node[inner sep=0pt] at (#1,0) {\Large$[$};}
\newcommand{\etiq}[2]{\node[below=10pt,inner sep=0pt] at (#1,0) {\scriptsize$#2$};}

\newpage

\pagestyle{fancy}

\fancyhf{} % clear existing header/footer entries
\fancyhead[R]{ {\Niveau}}
\fancyhead[L]{ {\Chapitre}}
\fancyfoot[L]{ A. Fernandez }
\fancyfoot[R]{{\thepage\  / \pageref{LastPage}}}


\section{La droite réelle}

\begin{definition}{Nombre réel}{}
  \begin{multihide}
  Un \textbf{nombre réel} est un nombre qui peut s'écrire avec une \textbf{partie entière} et une \textbf{partie décimale}. La \emph{partie décimale} peut contenir une liste \textbf{finie ou infinie} de chiffres.

  L'\textbf{ensemble} de tous les nombres réels est noté $\mathbb{R}$.

  Si un nombre $x$ est réel, on dit que \og $x$ appartient à $\mathbb{R}$ \fg{}, et on écrit \hide{$x \in \mathbb{R}$}.

  Sinon, on dit que \og $x$ n'appartient pas à $\mathbb{R}$ \fg{}, et on écrit \hide{$x \notin \mathbb{R}$}.
    
  \end{multihide}
\end{definition}

\begin{examples}
  \begin{itemize}
    \item[]
    \item \hideW{$0$, $3$, $-4$ sont des réels ayant une partie décimale nulle.}
    \item \hideW{$12{,}5$, $-1{,}33$, $0{,}4125$ sont des réels ayant une partie décimale qui contient un nombre fini de chiffres}
    \item \hideW{$1{,}2222\ldots$, $0{,}415415415\ldots$, $\dfrac{1}{3} \approx 0{,}333333\ldots$, $\sqrt{2} \approx 1{,}41421\ldots$ sont des réels ayant une partie décimale qui contient un nombre infini de chiffres.}
  \end{itemize}
\end{examples}

\begin{proposition}{}{}
  \begin{multihide}
  Tout nombre réel est l'abscisse d'un point d'une droite graduée appelée la \emph{droite numérique} ou \emph{droite réelle}.
  \end{multihide}
\end{proposition}

\begin{example}
  Placer sur la droite réelle, le plus précisément possible, les nombres :
  $$2{,}5 \qquad -1{,}5 \qquad -4{,}5 \qquad \dfrac{1}{3} \qquad \dfrac{10}{3} \qquad -\dfrac{1}{3} \qquad -\dfrac{12}{5} \qquad \sqrt{2}$$
\end{example}

\vspace{1em}

\begin{center}
  \begin{tikzpicture}[x=1.7cm]
    \draw[<->, >=Stealth, thick] (-5.4,0) -- (4.4,0);
    \foreach \i in {-5,...,4} {
      \draw (\i,0.1) -- (\i,-0.1);
      \node[below] at (\i,-0.1) {\small$\i$};
    }
    \hidetk{
      \foreach \x/\l in {2.5/2{,}5, -1.5/-1{,}5, -4.5/-4{,}5, 0.3333/\frac{1}{3}, 3.3333/\frac{10}{3}, -0.3333/-\frac{1}{3}, -2.4/-\frac{12}{5}, 1.41421/\sqrt{2}} {
        \fill (\x,0) circle[radius=2pt];
        \node[above=3pt] at (\x,0) {\small$\l$};
      }
    }
  \end{tikzpicture}
\end{center}

\section{Intervalles}

Un intervalle est un \og morceau \fg{} de la droite réelle contenant tous les nombres réels allant d'un nombre $a$ jusqu'à un nombre $b$. Par exemple, le morceau qui va de $-1$ à $5$ et qui contient tous les nombres réels allant de $-1$ à $5$ est un intervalle que l'on note à l'aide de crochets. L'orientation des crochets permet d'inclure ou non les extrémités $-1$ et $5$ de l'intervalle.

\begin{itemize}
  \item $[-1\,;\,5]$ est l'intervalle allant de $-1$ à $5$ dans lequel les nombres $-1$ et $5$ sont inclus.
  \item $]-1\,;\,5[$ est l'intervalle allant de $-1$ à $5$ dans lequel les nombres $-1$ et $5$ sont exclus.
  \item $[-1\,;\,5[$ est l'intervalle allant de $-1$ à $5$ dans lequel $-1$ est inclus et $5$ est exclu.
  \item $]-1\,;\,5]$ est l'intervalle allant de $-1$ à $5$ dans lequel $-1$ est exclu et $5$ est inclus.
\end{itemize}

Pour être plus précis, voici la définition rigoureuse des intervalles :

\newpage
\begin{definition}{Intervalle}{}
  Un \underline{intervalle} est un ensemble de nombres réels vérifiant une ou deux inégalités. Le tableau ci-dessous donne les différentes formes d'intervalles possibles, leur représentation et les inégalités qui les définissent. $a$ et $b$ sont deux réels tels que $a < b$.

  \vspace{-1em}
  \begingroup
  \renewcommand{\arraystretch}{1}
  \begin{center}
    \begin{tabular}{|C{2.6cm}|C{7cm}|C{3.4cm}|}
      \hline
      Intervalle & Représentation graphique & Inégalités \\ \hline
      $[a\,;\,b]$ &
      \begin{tikzpicture}
        \node at (0, 0.5) {$ $};
        \droite \draw[very thick] (1.5,0) -- (4.5,0);
        \brGF{1.5} \brDF{4.5} \etiq{1.5}{a} \etiq{4.5}{b}
      \end{tikzpicture}
      & \hideM{a \leq x \leq b} \\ \hline
      $]a\,;\,b[$ &
      \begin{tikzpicture}
        \node at (0, 0.5) {$ $};
        \droite \draw[very thick] (1.5,0) -- (4.5,0);
        \brGO{1.5} \brDO{4.5} \etiq{1.5}{a} \etiq{4.5}{b}
      \end{tikzpicture}
      & \hideM{a < x < b} \\ \hline
      $[a\,;\,b[$ &
      \begin{tikzpicture}
        \node at (0, 0.5) {$ $};
        \droite \hidetk{\draw[very thick] (1.5,0) -- (4.5,0);
        \brGF{1.5} \brDO{4.5}} \etiq{1.5}{a} \etiq{4.5}{b}
      \end{tikzpicture}
      & $a \leq x < b$ \\ \hline
      \hideM{]a\,;\,b]} &
      \begin{tikzpicture}
        \node at (0, 0.5) {$ $};
        \droite \draw[very thick] (1.5,0) -- (4.5,0);
        \brGO{1.5} \brDF{4.5} \etiq{1.5}{a} \etiq{4.5}{b}
      \end{tikzpicture}
      & $a < x \leq b$ \\ \hline
      $[a\,;\,+\infty[$ &
      \begin{tikzpicture}
        \node at (0, 0.5) {$ $};
        \droite \draw[very thick] (1.5,0) -- (5.9,0);
        \brGF{1.5} \etiq{1.5}{a}
      \end{tikzpicture}
      & \hideM{x \geq a} \\ \hline
      $]a\,;\,+\infty[$ &
      \begin{tikzpicture}
        \node at (0, 0.5) {$ $};
        \droite \hidetk{\draw[very thick] (1.5,0) -- (5.9,0);
        \brGO{1.5}} \etiq{1.5}{a}
      \end{tikzpicture}
      & $x > a$ \\ \hline
      \hideM{]-\infty\,;\,b]} &
      \begin{tikzpicture}
        \node at (0, 0.5) {$ $};
        \droite \draw[very thick] (0.1,0) -- (4.5,0);
        \brDF{4.5} \etiq{4.5}{b}
      \end{tikzpicture}
      & $x \leq b$ \\ \hline
      $]-\infty\,;\,b[$ &
      \begin{tikzpicture}
        \node at (0, 0.5) {$ $};
        \droite \hidetk{\draw[very thick] (0.1,0) -- (4.5,0);
        \brDO{4.5}} \etiq{4.5}{b}
      \end{tikzpicture}
      & $x < b$ \\ \hline
      \hideM{]-\infty\,;\,+\infty[} &
      \begin{tikzpicture}
        \node at (0, 0.5) {$ $};
        \droite \draw[very thick] (0.1,0) -- (5.9,0);
        \etiq{4.5}{}
      \end{tikzpicture}
      & \\ \hline
    \end{tabular}
  \end{center}
  \endgroup
\end{definition}

\begin{remarks}
  \begin{itemize}[noitemsep]
    \item[]
    \item À côté d'un signe $\infty$, le crochet est toujours ouvert.
    \item Si on veut inclure $a$ ou $b$, on met un crochet fermé du côté du nombre inclus.
    \item On a l'égalité d'ensembles : $\mathbb{R} = $ \hide{$]-\infty\,;\,+\infty[$}.
  \end{itemize}
\end{remarks}

\begin{definition}{Intersection et union}{}
  Soient $I$ et $J$ deux intervalles.
  \begin{itemize}[noitemsep]
    \item L'\underline{intersection} entre $I$ et $J$, notée $I \cap J$, est :
    
      \hide{l'ensemble des nombres réels qui appartiennent \textbf{à la fois} à $I$ et à $J$.}
    \item L'\underline{union} entre $I$ et $J$, notée $I \cup J$, est :
    
      \hide{l'ensemble des nombres réels qui appartiennent à $I$ \textbf{ou} à $J$ (au moins l'un des deux).}
  \end{itemize}
\end{definition}

\begin{exampleil}
  On se donne deux intervalles : $I = [1\,;\,7]$ et $J = [5\,;\,9]$.
  
  \vspace{1em}
  \begin{tikzpicture}[x=1.4cm]
    \draw[<->, >=Stealth, thick] (-1,0) -- (11.5,0);
    \foreach \i in {0,...,11} {
      \draw (\i,0.1) -- (\i,-0.1);
    }
    \node[below] at (0,-0.1) {$0$};
    % \hidetk{
    %   \draw[very thick, blue!70!black] (1,0.4) -- (7,0.4);
    %   \node[anchor=west, inner sep=0pt] at (1,0.4) {\Large$[$};
    %   \node[anchor=east, inner sep=0pt] at (7,0.4) {\Large$]$};
    %   \node[blue!70!black, above] at (4,0.5) {$I$};
    %   \draw[very thick, red!70!black] (5,0.9) -- (9,0.9);
    %   \node[anchor=west, inner sep=0pt] at (5,0.9) {\Large$[$};
    %   \node[anchor=east, inner sep=0pt] at (9,0.9) {\Large$]$};
    %   \node[red!70!black, above] at (7,1.0) {$J$};
    % }
  \end{tikzpicture}

\vspace{-1em} 
\begin{tasks}[label=\textbullet](2)
  \task $I \cap J = \hideM{[5\,;\,7].}$
  \task $I \cup J = \hideM{[1\,;\,9].}$
\end{tasks}
\end{exampleil}

\begin{remarks}
  \begin{itemize}[noitemsep]
    \item[]
    \item Si deux intervalles n'ont aucun nombre en commun, leur intersection est \hide{l'\textbf{ensemble vide}}, noté \hide{$\varnothing$}.

      Par exemple : $[1\,;\,3] \cap [5\,;\,8] = \hideM{\varnothing}$ et $]-\infty\,;\,2[\, \cap\, [2\,;\,+\infty[\, = \hideM{\varnothing}$.
    \item L'union de deux intervalles n'est pas toujours un intervalle. Par exemple, $[1\,;\,3] \cup [5\,;\,8]$ \hide{n'est pas un intervalle} : il y a un \og trou \fg{} entre $3$ et $5$.
  \end{itemize}
\end{remarks}

\subsection*{Inégalités et opérations}

Pour savoir à quel intervalle appartient une expression comme $2x+1$ lorsque $x$ appartient à un intervalle, on transforme les inégalités à l'aide de la propriété suivante.

\begin{proposition}{Inégalités et opérations}{}
  Soient $a$, $b$ et $c$ trois nombres réels tels que $a \leq b$.
  \begin{itemize}
    \item \textbf{Ajouter ou soustraire} un même nombre \hide{ne change pas le sens de l'inégalité :}
    $$\hideM{a + c \leq b + c \text{ et } a - c \leq b - c}$$
    \item \textbf{Multiplier ou diviser} par un même nombre \textbf{strictement positif} \hide{ne change pas le sens de l'inégalité :}
    $$\hideM{\text{si } c > 0, \text{ alors } a \times c \leq b \times c}$$
    \item \textbf{Multiplier ou diviser} par un même nombre \textbf{strictement négatif} \hide{\underline{change le sens} de l'inégalité :}
    $$\hideM{\text{si } c < 0, \text{ alors } a \times c \geq b \times c}$$
  \end{itemize}
\end{proposition}

\begin{example}
  On sait que $x \in [1\,;\,3]$. À quel intervalle appartient $-2x + 5$ ?
  
  \vspace{0.2em}
  \framebox[1\linewidth]{\parbox{5.5in}{%\centering \underline{Réponse:}
  \[
    \hideMW{1 \leq x \leq 3
    \quad \Leftrightarrow \quad -2 \geq -2x \geq -6
    \quad \Leftrightarrow \quad 3 \geq -2x + 5 \geq -1}
  \]
  \hideW{On a multiplié par $-2$ : le sens change. Ainsi $-2x + 5 \in [-1\,;\,3]$}.
  %\vspace{-0.5em}
  }}
\end{example}


% \newpage
% 
% \section{Valeur absolue}
% 
% \begin{definition}{Valeur absolue et distance}{}
%   \begin{multihide}
%   \begin{itemize}
%     \item Soit $a$ un nombre réel. La \underline{valeur absolue} de $a$, notée $|a|$, est la longueur qui sépare $a$ de $0$. Autrement dit, $|a|$ est simplement la valeur positive de $a$.
%     \item Soient $a$ et $b$ deux nombres réels. On appelle \underline{distance} entre $a$ et $b$ la valeur $|a - b|$.
%   \end{itemize}
%   \end{multihide}
% \end{definition}
% 
% \begin{example}
%   On a placé deux réels $a$ et $b$ sur la droite réelle. On a choisi $a = -5$ et $b = 2$.
% 
% \begin{center}
%   \begin{tikzpicture}[x=1.3cm]
%     \draw[<->, >=Stealth, thick] (-6.4,0) -- (4.4,0);
%     \foreach \i in {-6,...,4} {
%       \draw (\i,0.1) -- (\i,-0.1);
%     }
%     \node[below] at (-5,-0.1) {$a = -5$};
%     \node[below] at (0,-0.1) {$0$};
%     \node[below] at (2,-0.1) {$b = 2$};
%   \end{tikzpicture}
% \end{center}
% \vspace{0.2em}
% \framebox[1\linewidth]{\parbox{5.5in}{%\centering \underline{Réponse:}
% \begin{multihideW}
% \begin{itemize}
%   \item La distance entre $a$ et $0$ vaut $5$. Ainsi : $|a| = |$\hide{$-5$}$| = $\hide{$5$}.
%   \item La distance entre $a$ et $b$ vaut $7$. Ainsi : $|a - b| = |$\hide{$-5 - 2$}$| = |$\hide{$-7$}$| = $\hide{$7$}.
% \end{itemize}
% \end{multihideW}
% }}
% 
% \end{example}
% 
% \begin{remarks}
%   \begin{itemize}
%     \item[]
%     \item Soit $a$ un réel. Si $a$ est positif alors $|a| = a$. Si $a$ est négatif alors $|a| = -a$.
% 
%     La valeur absolue est donc toujours une valeur positive.
%     \item Soient $a$ et $b$ deux réels. Alors $|a - b| = |b - a|$.
% 
%     C'est cohérent, la distance de $a$ à $b$ est bien la même que celle de $b$ à $a$.
%   \end{itemize}
% \end{remarks}
% 
% \begin{proposition}{}{}
%   \begin{multihide}
%   Soient $a$ et $x$ deux réels et $r$ un réel positif.
%   $$x \in [a - r\,;\,a + r] \quad \text{si et seulement si} \quad |x - a| \leq r$$
%     
%   \end{multihide}
% \end{proposition}
% 
% \begin{remark}
%   \og Si et seulement si \fg{} a une signification très précise en mathématiques. Son utilisation signifie que les deux affirmations situées avant et après sont exactement équivalentes. Autrement dit, la propriété ci-dessus stipule à la fois que :
%   \begin{center}
%     Si $|x - a| \leq r$ alors $x \in [a - r\,;\,a + r]$
%   \end{center}
%   et que réciproquement :
%   \begin{center}
%     Si $x \in [a - r\,;\,a + r]$ alors $|x - a| \leq r$
%   \end{center}
%   Si la propriété était : \og $x \in [a - r\,;\,a + r]$ si $|x - a| \leq r$ \fg{}, celle-ci affirmerait que le premier point est vrai, mais pas forcément la réciproque.
% \end{remark}
% 
% \begin{example}
%   On a représenté la droite réelle ci-dessous :
% 
% \begin{center}
%   \begin{tikzpicture}[x=1.3cm]
%     \draw[<->, >=Stealth, thick] (-4.4,0) -- (6.4,0);
%     \foreach \i in {-4,...,6} {
%       \draw (\i,0.1) -- (\i,-0.1);
%     }
%     \node[below] at (-2,-0.1) {$a - r$};
%     \node[below] at (0,-0.1) {$0$};
%     \node[below] at (1,-0.1) {$a$};
%     \node[below] at (4,-0.1) {$a + r$};
%   \end{tikzpicture}
% \end{center}
% 
% \framebox[1\linewidth]{\parbox{5.5in}{%\centering \underline{Réponse:}
% \begin{multihideW}
% Ici, on a : $a = 1$ et $r = 3$.
% 
% On voit bien qu'il est équivalent de dire qu'un nombre $x$ appartient à l'intervalle $[a-3\,;\,a+3]$ ou de dire que la distance entre $x$ et $a$ est plus petite que $3$.
% \end{multihideW}
% }}
% \end{example}
% 
% \vfill

\newpage

\section{Les sous-ensembles particuliers de $\mathbb{R}$}

\begin{definition}{Sous-ensemble de $\mathbb{R}$}{}
  \begin{multihide}
  Un \emph{sous-ensemble} de $\mathbb{R}$ est un ensemble inclus dans $\mathbb{R}$. Autrement dit, c'est un ensemble constitué de nombres réels (et qui ne les contient en général pas tous).
  \end{multihide}
\end{definition}

\begin{remark}
  Les intervalles sont des exemples de sous-ensembles de $\mathbb{R}$.
\end{remark}

\begin{definition}{Principaux sous-ensembles de $\mathbb{R}$}{}
  On définit ci-dessous les principaux sous-ensembles de $\mathbb{R}$ à connaître.
  \begin{itemize}
    \item L'ensemble des \textbf{entiers naturels}, noté $\mathbb{N}$, est constitué de tous les nombres entiers positifs ou nuls :
    $$\hideM{\mathbb{N} = \{0\,;\,1\,;\,2\,;\,3\,;\,\ldots\}}$$
    \item L'ensemble des \textbf{entiers relatifs}, noté $\mathbb{Z}$, est constitué de tous les nombres entiers positifs, nuls ou négatifs :
    $$\hideM{\mathbb{Z} = \{\ldots\,;\,-3\,;\,-2\,;\,-1\,;\,0\,;\,1\,;\,2\,;\,3\,;\,\ldots\}}$$
    \item L'ensemble des \textbf{décimaux}, noté $\mathbb{D}$, est constitué de tous les nombres réels pouvant s'écrire sous la forme \hide{$\dfrac{a}{10^n}$ où $a$ est un entier relatif et $n$ un entier naturel :}
    $$\hideM{\mathbb{D} = \left\{\dfrac{a}{10^n} \text{ où } a \in \mathbb{Z},\ n \in \mathbb{N}\right\}}$$
    \item L'ensemble des \textbf{rationnels}, noté $\mathbb{Q}$, est constitué de tous les nombres réels pouvant s'écrire sous la forme \hide{d'une fraction $\dfrac{a}{b}$ où $a$ est un entier relatif et $b$ un entier naturel différent de $0$ :}
    $$\hideM{\mathbb{Q} = \left\{\dfrac{a}{b} \text{ où } a \in \mathbb{Z},\ b \in \mathbb{N}^*\right\}}$$
  \end{itemize}
\end{definition}

\begin{examples}
  \begin{tasks}[label=\textbullet](2)
    \task $17{,}2\hideM{ \in \mathbb{D} \text{ car } 17{,}2 = \dfrac{172}{10}}$
    \task $-0{,}333\hideM{ \in \mathbb{D} \text{ car } -0{,}333 = \dfrac{-333}{10^3}}$
    \task $\dfrac{1}{2}\hideM{ \in \mathbb{Q}}$
    \task $\dfrac{1}{2}\hideM{ \in \mathbb{D} \text{ car } \dfrac{1}{2} = \dfrac{5}{10}}$
    \task $\dfrac{21}{3}\hideM{ \in \mathbb{Q}}$
    \task $\dfrac{21}{3}\hideM{ \in \mathbb{N} \text{ car } \dfrac{21}{3} = 7}$
  \end{tasks}
\end{examples}

\begin{remark}
  Les nombres décimaux ne sont pas seulement ceux qui s'écrivent avec une partie décimale ayant un nombre fini de chiffres. Par exemple $0{,}999999\ldots$ est un nombre décimal bien qu'il ait une partie décimale avec une infinité de $9$ puisque $0{,}999999\ldots = 1$.

  En revanche, il est vrai que tout nombre décimal peut s'écrire avec une partie décimale ayant un nombre fini de chiffres.
\end{remark}

\begin{proposition}{}{}
  \begin{multihide}
  $$\mathbb{N} \subset \mathbb{Z} \subset \mathbb{D} \subset \mathbb{Q} \subset \mathbb{R}$$
  où le symbole $\subset$ signifie \og est inclus dans \fg{}.
  \end{multihide}
\end{proposition}

\begin{remarks}
  \begin{itemize}[noitemsep]
    \item[]
    \item Le schéma ci-dessous représente les inclusions des différents sous-ensembles de $\mathbb{R}$, les uns dans les autres. Cela montre que les entiers naturels font partie des entiers relatifs, qui font eux-mêmes partie des décimaux, faisant eux-mêmes partie des rationnels, faisant eux-mêmes partie des réels.

\begin{center}
  \begin{tikzpicture}[scale=1]
    % Ensembles emboîtés
    \draw[blue!60!black, thick, rounded corners=12pt] (0,0) rectangle (11.5,7);
    \draw[magenta!70!black, thick, rounded corners=12pt] (0.3,0.3) rectangle (9.6,6);
    \draw[magenta!70!black, thick, rounded corners=10pt] (0.6,0.6) rectangle (7.6,4.9);
    \draw[magenta!70!black, thick, rounded corners=8pt] (0.9,0.9) rectangle (6,3.6);
    \draw[magenta!70!black, thick, rounded corners=6pt] (1.2,1.2) rectangle (3.2,2.7);
    % Étiquettes
    \node[fill=white, inner sep=2pt] at (1.5,7)   {\textcolor{blue!60!black}{$\mathbb{R}$}};
    \node[fill=white, inner sep=2pt] at (1.5,6)   {\textcolor{magenta!70!black}{$\mathbb{Q}$}};
    \node[fill=white, inner sep=2pt] at (1.5,4.9) {\textcolor{magenta!70!black}{$\mathbb{D}$}};
    \node[fill=white, inner sep=2pt] at (1.5,3.6) {\textcolor{magenta!70!black}{$\mathbb{Z}$}};
    \node[fill=white, inner sep=2pt] at (1.7,2.7) {\textcolor{magenta!70!black}{$\mathbb{N}$}};
    % Éléments de N
    \node at (1.7,1.7) {\hide{$0$}};
    \node at (2.6,2.2) {\hide{$15$}};
    \node at (2.5,1.5) {\hide{$2^3$}};
    % Éléments de Z
    \node at (1.6,3.1) {\hide{$-1$}};
    \node at (3.6,3.1) {\hide{$-12$}};
    \node at (4.7,1.8) {\hide{$-3^4$}};
    % Éléments de D
    \node at (2.4,4.3) {\hide{$-0{,}475$}};
    \node at (5.5,4.3) {\hide{$10^{-3}$}};
    \node at (6.8,2.4) {\hide{$-\frac{3}{4}$}};
    % Éléments de Q
    \node at (8.6,5.2) {\hide{$\frac{22}{7}$}};
    \node at (8.6,2.4) {\hide{$-\frac{4}{3}$}};
    % Éléments de R
    \node at (4.5,6.5) {\hide{$\sqrt{2}$}};
    \node at (8,6.5)    {\hide{$\cos 30^\circ$}};
    \node at (10.6,5.2) {\hide{$\pi$}};
    \node at (10.6,1.5) {\hide{$-\sqrt{\frac{3}{4}}$}};
  \end{tikzpicture}
\end{center}
  \item Le fait que tout nombre décimal puisse s'écrire sous la forme $\dfrac{a}{10^n}$ prouve que tout nombre décimal est un nombre rationnel. On a donc bien $\mathbb{D} \subset \mathbb{Q}$. L'inverse est-il vrai ? Tout nombre rationnel est-il un nombre décimal ? Autrement dit, a-t-on $\mathbb{Q} \subset \mathbb{D}$ ? La propriété suivante répond à cette question.
\end{itemize}
\end{remarks}

\begin{proposition}{}{}
  Le nombre $\dfrac{1}{3}$ n'est pas décimal.
\end{proposition}

%\needspace{4\baselineskip}
\begin{proof}$ $\newline
\vspace{0.2em}
\noindent\framebox[\linewidth]{\parbox{\dimexpr\linewidth-6\fboxsep\relax}{
\begin{multihideW}
  Supposons que $\frac{1}{3}$ soit décimal. Il existerait alors un entier $a$ et un entier naturel $n$ tels que
  $$\dfrac{1}{3} = \dfrac{a}{10^n}.$$
  Par un produit en croix, on obtient $10^n = 3a$ : le nombre $10^n$ serait donc un multiple de $3$.

  Or $10^n$ s'écrit avec un $1$ suivi de $n$ zéros : la somme de ses chiffres vaut $1$, qui n'est pas un multiple de $3$. D'après le critère de divisibilité par $3$, $10^n$ n'est donc pas un multiple de $3$.

  C'est une contradiction. La supposition de départ est donc fausse : $\frac{1}{3}$ n'est pas décimal.
\end{multihideW}
}}
\end{proof}

\begin{remark}
  On utilise un \textbf{raisonnement par l'absurde} : pour démontrer qu'une affirmation est vraie, on suppose au contraire qu'elle est fausse, puis on en déduit quelque chose d'impossible (une \emph{contradiction}). Puisqu'on aboutit à une impossibilité, c'est que la supposition de départ était fausse : l'affirmation est donc vraie.
\end{remark}

\begin{remark}
  Cela prouve qu'il existe un nombre rationnel qui n'est pas décimal et donc que $\mathbb{Q} \not\subset \mathbb{D}$.

  À l'Antiquité, les Grecs ont pensé pendant des siècles que tout nombre pouvait s'écrire comme une fraction entre deux entiers. En termes modernes, ils pensaient qu'on avait l'égalité d'ensembles $\mathbb{R} = \mathbb{Q}$. C'est vers 530 avant J.-C. que l'un d'entre eux réussit à prouver que la longueur de la diagonale d'un carré de côté $1$ ne pouvait pas s'écrire comme une fraction. Autrement dit, il prouva la propriété suivante :
\end{remark}

\begin{proposition}{(Admise)}{}
  Le nombre $\sqrt{2}$ n'est pas un nombre rationnel.
\end{proposition}

\begin{remark}
  Cela prouve qu'il existe un nombre réel qui n'est pas rationnel et donc que $\mathbb{R} \not\subset \mathbb{Q}$. Les nombres réels qui ne sont pas rationnels se nomment les \emph{irrationnels}. Les nombres suivants sont d'autres exemples de nombres irrationnels : $\sqrt{3}$, $\sqrt{5}$, $\pi$.

  La démonstration de cette propriété sera faite en fin d'année, dans le chapitre d'arithmétique.
\end{remark}

\section{Encadrement décimal d'un nombre réel}

Un nombre réel comme $\dfrac{1}{3}$ ou $\sqrt{2}$ a une infinité de chiffres après la virgule : on ne peut pas l'écrire exactement sous forme décimale. On peut en revanche le \og coincer \fg{} entre deux nombres décimaux aussi proches que l'on veut.

\begin{definition}{Encadrement décimal}{}
  \begin{multihide}
  Soient $x$ un nombre réel et $n$ un entier naturel.

  Un \textbf{encadrement décimal de $x$ à $10^{-n}$ près} (on dit aussi \textbf{d'amplitude $10^{-n}$}) est une double inégalité
  $$a \leq x \leq b$$
  où $a$ et $b$ sont deux nombres décimaux ayant $n$ chiffres après la virgule et tels que $b - a = 10^{-n}$.
  \end{multihide}
\end{definition}

\begin{remark}
  L'\textbf{amplitude} d'un encadrement $a \leq x \leq b$ est la différence $b - a$. Plus l'amplitude est petite, plus l'encadrement est précis.

  $10^{-1} = 0{,}1$ (au dixième près), $10^{-2} = 0{,}01$ (au centième près), $10^{-3} = 0{,}001$ (au millième près).
\end{remark}

\begin{method}{Encadrer un nombre à $10^{-n}$ près}{}
  \begin{multihide}
  \begin{itemize}
    \item On écrit le nombre avec \textbf{au moins} $n + 1$ chiffres après la virgule (en posant la division pour une fraction).
    \item On \textbf{coupe} après le $n$-ième chiffre après la virgule : on obtient une borne de l'encadrement.
    \item On obtient l'autre borne en ajoutant ou en retranchant $10^{-n}$, de façon que le nombre soit bien compris entre les deux.
  \end{itemize}
  \end{multihide}
\end{method}

\begin{examples}
  \begin{itemize}
    \item[]
    \item Encadrement de $12{,}4579$ à $10^{-2}$ près : \hideM{12{,}45 \leq 12{,}4579 \leq 12{,}46}
    \item Encadrement de $\dfrac{2}{3}$ à $10^{-3}$ près. En posant la division, $\dfrac{2}{3} = 0{,}6666\ldots$ donc \hideM{0{,}666 \leq \dfrac{2}{3} \leq 0{,}667}
    \item Encadrement de $-5{,}783$ à $10^{-1}$ près : \hideM{-5{,}8 \leq -5{,}783 \leq -5{,}7}

      \hide{Attention : pour un nombre négatif, couper après la virgule donne la borne de \textbf{droite}.}
  \end{itemize}
\end{examples}

\vfill

\begin{definition}{Arrondi}{}
  \begin{multihide}
  L'\textbf{arrondi} de $x$ à $10^{-n}$ près est, parmi les deux bornes de l'encadrement de $x$ à $10^{-n}$ près, celle qui est \textbf{la plus proche} de $x$. Si $x$ est exactement au milieu, on choisit la borne de droite.
  \end{multihide}
\end{definition}

\begin{example}
  $12{,}45 \leq 12{,}4579 \leq 12{,}46$ et $12{,}4579$ est plus proche de $12{,}46$ : l'arrondi de $12{,}4579$ au centième est \hide{$12{,}46$}.
\end{example}

\vspace{2em}

\begin{savoir}{}{}
  \begin{itemize}
    \item Placer des nombres réels sur la droite réelle.
    \item Représenter un intervalle et traduire un intervalle par des inégalités (et réciproquement).
    \item Déterminer l'intersection et l'union de deux intervalles.
    \item Transformer une inégalité en ajoutant ou en multipliant par un nombre (attention au signe).
    \item Connaître les ensembles $\mathbb{N}$, $\mathbb{Z}$, $\mathbb{D}$, $\mathbb{Q}$, $\mathbb{R}$ et leurs inclusions, et savoir dans quel(s) ensemble(s) se trouve un nombre donné.
    \item Donner un encadrement décimal d'un nombre réel à $10^{-n}$ près, et son arrondi.
  \end{itemize}
\end{savoir}

\vfill

\end{document}